\section{结语：华北统治是一项持续的接收}

金的国家能力来自把征服转为中都的供给、地方税役、军政组织和边界防务，而非来自一次军事胜利本身。与南宋的和议和战争把淮河、长江、华北道路与东南财赋放在同一竞争场中；东北、山东、河南和中都周边又不以同一种方式进入国家账目。因而，金既不能只被写成南宋的敌手，也不能只被写作草原政治的延长。

蒙古战争暴露的不是一个抽象“北方王朝”的必然衰败，而是税源、运输和动员在多线压力下失去协调的具体过程。金亡之后，城市、军户、工匠和地方社会并未同时消失；他们进入蒙古接收与南宋防务的不同路径。下一章将从南宋的江海财政与边防重新审视这一互动，而不是把南方留作金史的余景。
